\chapter{Vektorindexierung mit OpenSearch}
\label{ch:indexierung}
% Deckt Ziel 3 ab: "Indexierung der ausgewaehlten Teilmenge als dichte
% Vektoren mittels OpenSearch".

Damit die im vorherigen Pipeline-Schritt ausgewählte Wikipedia-Teilmenge von einem Large-Language-Modell genutzt
werden kann, müssen diese Artikel zunächst in geeignete Textabschnitte (Chunks) unterteilt werden, bevor diese
als dichte Vektoren indexiert werden können. Die Aufteilung ermöglicht eine granularere Suche nach relevanten
Informationen, da bei einer Anfrage nicht zwangsläufig ein vollständiger Wikipedia-Artikel, sondern gezielt
relevante Textpassagen als Kontext für das Large-Language-Modell abgerufen werden können. Die Verwendung von
Textpassagen als Retrieval-Einheiten stellt dabei einen etablierten Ansatz im Dense Retrieval dar
~\cite{karpukhin-etal-2020-dense}.

\section{Chunking-Strategie}
\label{sec:chunking-strategie}

Die Zerlegung der Wikipedia-Artikel erfolgt abschnittsweise. Die Methode
\texttt{chunk\_article()} verarbeitet dazu die im Artikelobjekt gespeicherten
Abschnitte einzeln. Dadurch wird
sichergestellt, dass ein Chunk ausschließlich Text aus einem einzelnen Wikipedia-Abschnitt enthält und
der inhaltliche Zusammenhang der jeweiligen Passage erhalten bleibt.

Innerhalb eines Abschnitts werden die Texte mit einer maximalen Fenstergröße von 250 Wörtern und einer
Überlappung von 40 Wörtern unterteilt. Die Wahl der Chunk-Größe stellt einen Kompromiss zwischen der
Granularität des Retrievals und dem Erhalt ausreichenden Kontextes dar. Untersuchungen zeigen, dass die
optimale Chunk-Größe von der Charakteristik der Daten, der konkreten Retrieval-Aufgabe und dem verwendeten
Embedding-Modell abhängt und daher nicht allgemein festgelegt werden kann
~\cite{bhat2025rethinkingchunksizelongdocument}. Die für dieses Projekt verwendete Fenstergröße wurde
daher als geeigneter Kompromiss für den verwendeten Wikipedia-Korpus und das nachfolgende Retrieval
festgelegt.

Die Überlappung soll verhindern, dass zusammengehörige Informationen durch eine Chunk-Grenze vollständig
voneinander getrennt werden. Zusätzlich enthält jeder Chunk Metadaten, insbesondere den Titel des
zugehörigen Wikipedia-Abschnitts. Diese werden im weiteren Verlauf der Pipeline verwendet, um die Herkunft
der abgerufenen Textpassagen im Prompt kenntlich zu machen und eine Quellenattribution zu ermöglichen.


\section{Embedding-Modell}
\label{sec:embedding-modell}

Als Embedding-Modell kommt \texttt{sentence-transformers/all-MiniLM-L6-v2} zum Einsatz, welches jeden
Chunk-Text auf einen 384-dimensionalen dichten Vektor abbildet~\cite{minilm-model-card}. Mit 22,7 Millionen
Parametern ist es deutlich kleiner als vergleichbare Modelle wie \texttt{all-mpnet-base-v2}. Die offizielle
Dokumentation der \texttt{sentence-transformers}-Bibliothek beschreibt \texttt{all-MiniLM-L6-v2} als
schnellere Alternative zu \texttt{all-mpnet-base-v2}, während letzteres eine höhere Qualität bei entsprechend
höherem Rechenaufwand bietet~\cite{sbert-pretrained-models}.

Für das Projekt ist insbesondere der geringere Rechenaufwand relevant, da sowohl das initiale Embedding der
ausgewählten Chunks während der Indexierung als auch das Embedding eingehender Anfragen auf der vorgesehenen
CPU-only-Hardware ausgeführt werden muss. Das Modell \texttt{all-MiniLM-L6-v2} stellt hierfür einen geeigneten
Kompromiss zwischen Modellgröße und Embedding-Qualität dar.

Das Modell wurde auf über 1,17~Milliarden Satzpaaren aus unterschiedlichen Quellen trainiert und ist unter
anderem für Information-Retrieval-, Clustering- und Sentence-Similarity-Aufgaben vorgesehen
~\cite{minilm-model-card}. Damit ist es für den im Projekt verfolgten Anwendungsfall des semantischen
Retrievals geeignet.

Eine weitere relevante Einschränkung des Modells ist die maximale Eingabelänge von 256 Word Pieces
~\cite{minilm-model-card}. Da die Chunk-Größe anhand von Wörtern festgelegt wird, kann daraus jedoch nicht
direkt abgeleitet werden, dass jeder Chunk diese Grenze unterschreitet, da ein Wort in mehrere Word Pieces
zerlegt werden kann. Die tatsächliche Eingabelänge wird daher durch das verwendete Tokenisierungsverfahren
begrenzt.

Für die Indexierung und die spätere Abfrage wird dasselbe Embedding-Modell mit identischen Einstellungen
verwendet. Sowohl \texttt{embed\_chunks()} als auch \texttt{embed\_query()} greifen auf dieselbe
\texttt{Embedder}-Instanz mit aktivierter L2-Normalisierung
(\texttt{normalize\_embeddings=True}) zurück. Dadurch werden Query- und Chunk-Vektoren im selben Vektorraum
repräsentiert. Bei L2-normalisierten Vektoren entspricht das Skalarprodukt der Kosinus-Ähnlichkeit, wodurch
die Berechnung direkt mit dem im OpenSearch-Index verwendeten \texttt{cosinesimil}-Ähnlichkeitsmaß
übereinstimmt (Abschnitt~\ref{sec:opensearch-index-design}).

Zur Einschätzung des Rechenaufwands wurde die Embedding-Geschwindigkeit isoliert gemessen -- also ohne
den Anteil von OpenSearch-Bulk-Writes oder ZIM-Zugriff, um ausschließlich die Kosten der Modellinferenz
abzubilden. Unterschieden wurden zwei Fälle: der Massen-Durchsatz beim Indexieren (Batch-Embedding mit
\texttt{embed\_chunks()}, Chunks von ca.\ 250 Wörtern Länge entsprechend der in
Abschnitt~\ref{sec:chunking-strategie} festgelegten Fenstergröße) sowie die Latenz einer einzelnen
Anfrage (\texttt{embed\_query()}, Batch-Größe 1, wie im Retrieval-Pfad zur Laufzeit).

Die Messung erfolgte auf dem für die Indexierung des vollen Korpus verwendeten Build-Server (104~Kerne),
nicht auf dem Raspberry Pi selbst -- die Zahlen belegen daher die grundsätzliche Eignung des Modells,
nicht die konkrete Laufzeit auf der Zielhardware. Beim Batch-Embedding wurden 655,0~Chunks pro Sekunde
erreicht, bei einzelnen Anfragen betrug die mittlere Latenz 5,4~ms pro Query. Insbesondere die
Query-Latenz liegt damit deutlich unterhalb der in Kapitel~\ref{ch:evaluation} gemessenen
Gesamt-Retrieval-Latenz, sodass das Embedding selbst auf dieser Hardware keinen dominanten Anteil an der
Anfragebearbeitung ausmacht. Eine entsprechende Messung direkt auf dem Raspberry Pi steht noch aus.


\section{OpenSearch-Index-Design}
\label{sec:opensearch-index-design}
Für die Vektorsuche wird ein Feld vom Typ \texttt{knn\_vector} mit der Methode \texttt{hnsw} verwendet.
Der Grund für ein approximatives statt eines exakten k-NN-Verfahrens liegt in der Skalierbarkeit: Laut der OpenSearch-Dokumentation 
berechnet exaktes k-NN die Distanz per Brute-Force zwischen der Anfrage und allen Punkten, was bei 
\enquote{extremely large datasets with high dimensionality \ldots\ a scaling problem that reduces the efficiency of the search} erzeugt. 
Approximative Verfahren \enquote{overcome this by employing tools that restructure indexes more efficiently 
\ldots\ requir[ing] a sacrifice in accuracy but increas[ing] search processing speeds appreciably}~\cite{opensearch-approximate-knn}.

Dieser Trade-off zulasten der Genauigkeit ist in unserem Anwendungsfall jedoch sinnvoll, da laut Dokumentation approximatives 
k-NN \enquote{the preferred method \ldots\ when a dataset reaches hundreds of thousands of vectors} darstellt~\cite{opensearch-approximate-knn}.

Als \texttt{space\_type} wird \texttt{cosinesimil} verwendet, definiert als 
$d(x,y) = 1-\cos\theta = 1 - \frac{x\cdot y}{|x|\cdot|y|}$ ~\cite{opensearch-spaces}. Da die im Embedding-Schritt (Abschnitt~\ref{sec:embedding-modell}) 
erzeugten Vektoren bereits L2-normalisiert sind ($|x|=|y|=1$), reduziert sich diese Distanzfunktion auf ein einfaches Skalarprodukt

Die HNSW-Parameter \texttt{ef\_construction} und \texttt{m} werden im Mapping nicht explizit gesetzt und laufen damit auf den \texttt{nmslib}-Standardwerten
$\texttt{ef\_construction}=100$, $m=16$ ~\cite{opensearch-methods-engines}. Höhere Werte würden laut Dokumentation einen genaueren Graphen bei langsamerer 
Indexierungsgeschwindigkeit (\texttt{ef\_construction}) bzw. höherem Speicherverbrauch ($m$) erkaufen. 

Für den Bulk-Indexiervorgang (\texttt{index\_chunks()}) wird die deterministische
\texttt{chunk\_id} jedes Chunks als Dokument-ID in OpenSearch verwendet.
Ein wiederholter Indexierlauf überschreibt so bestehende Dokumente, statt
Duplikate anzulegen, was praktisch für wiederholte Testläufe ist. \texttt{create\_index()} legt den Index nur an, wenn er noch nicht existiert; ein \texttt{recreate}-Flag
(CLI: \texttt{--recreate-index}) erlaubt bei Bedarf das gezielte Neuanlegen, etwa nach einer Mapping-Änderung.

Der lokale \texttt{docker-compose.yml}-Aufbau, mit dem während der Entwicklung getestet wurde, begrenzt den OpenSearch-Java-Heap auf 2~GB — für den 
kleinen Testkorpus auf dem eigenen Rechner ausreichend. Der eigentliche Indexierungslauf über den vollen Korpus fand aber nicht lokal und nicht auf 
dem Pi statt, sondern auf einem für diesen Zweck bereitgestellten leistungsstärkeren Server, wodurch Speicherknappheit als limitierender Faktor
eliminiert wurde. Der fertige Index wurde anschließend per OpenSearch-Snapshot vom Server auf den Pi übertragen (\texttt{path.repo}-Mechanismus), 
statt den aufwendigen Indexierungsvorgang selbst auf der eingeschränkten Zielhardware laufen zu lassen.

\section{Skalierung auf den vollen Korpus}
Dieser Abschnitt betrifft nicht mehr den vollen, ungefilterten Wikipedia-Korpus, sondern die bereits durch die Selection (Kapitel~\ref{ch:dokumentenauswahl}) 
auf das konfigurierte Speicherbudget reduzierte Teilmenge. Die aufwendige Zwei-Pass-Streaming-Architektur, die dort nötig war, greift hier nicht mehr, denn 
Chunking, Embedding und Indexierung laufen über eine von vornherein kleinere, budgetbeschränkte Artikelmenge.

Im \texttt{index}-Kommando der CLI wird die Chunk-Liste für diese gesamte ausgewählte Teilmenge auf einmal aufgebaut, statt gestreamt zu werden. Erst 
Embedding und Indexierung selbst laufen anschließend batchweise (in Schritten von \texttt{config.embeddings.batch\_size}, standardmäßig 32): pro 
Durchlauf werden nur die Vektoren eines Batches berechnet und indexiert, nie die aller Chunks gleichzeitig.

Dass die vollständig im Speicher gehaltene Chunk-Liste hier kein Problem darstellte, hat vor allem zwei Gründe: erstens ist die Teilmenge 
durch das Speicherbudget von Haus aus begrenzt, zweitens lief der eigentliche Indexierungslauf nicht auf dem Pi, sondern auf einem dafür 
bereitgestellten leistungsstärkeren Server.

\section{Retrieval}
% Query-Embedding + k-NN-Suche, Retriever-Klasse als Abstraktionsschicht.